% ============================================================
% Paper 2 — Causal B-JEPA
% Uncertainty-Aware Gene Regulatory Network Inference
% from Interventional and Observational Data
%
% Timeline:
%   Summer 2026   — Build architecture, run experiments
%   Oct    2026   — NeurIPS MLCB preliminary results (if ready)
%   Early  2027   — Genome Biology / Nature Computational Science / NeurIPS 2027
%
% Target length: ~15 pages main text + supplementary
% ============================================================

\documentclass[11pt]{article}
\usepackage[margin=1in]{geometry}
\usepackage{amsmath, amssymb, amsthm}
\usepackage{graphicx}
\usepackage{hyperref}
\usepackage{booktabs}
\usepackage{natbib}
\usepackage{xcolor}
\usepackage{algorithm}
\usepackage{algorithmic}

\newcommand{\todo}[1]{\textcolor{red}{[TODO: #1]}}
\newcommand{\placeholder}[1]{\textcolor{gray}{[PLACEHOLDER: #1]}}

\title{Causal B-JEPA: Sparse Attention World Models for\\
       Uncertainty-Quantified Gene Regulatory Network Inference}
\author{Raja Giddi}
\date{2027}

\begin{document}
\maketitle

% ============================================================
\begin{abstract}
% Central claim (4 sentences):
% 1. Standard world models fail for GRN inference because they implement
%    dense global prediction when the problem requires sparse local causal identification.
% 2. Causal B-JEPA addresses three separable failure modes from Paper 1 simultaneously.
% 3. Sparsemax attention weights are the regulatory network; NUTS on the sparse
%    support produces calibrated credible intervals — the first in GRN inference.
% 4. Comparison of observational IV and interventional training quantifies the
%    causal identification gap for specific edge classes.
% ~250 words.
\todo{Abstract}
\end{abstract}

% ============================================================
\section{Introduction}
% Paragraph 1: Paper 1 finding — dense world models fail for sparse causal inference.
%   This paper builds the architecture Paper 1 specified.
%   One sentence establishing prior claim: cite Paper 1.
%
% Paragraph 2: The three components GRN inference requires.
%   Causal identification strategy (IV or interventional).
%   Sparse local output head in expression space.
%   Tractable calibrated posteriors over sparse support.
%
% Paragraph 3: Causal B-JEPA delivers all three.
%   Two versions (A: observational IV, B: interventional).
%   First calibrated credible intervals over regulatory edges.
%   Uncertainty-guided experimental prioritization.
%
% Paragraph 4: Paper contributions enumerated.
\todo{Introduction}

% ============================================================
\section{Background}

\subsection{Failure Analysis of Dense World Models for GRN Inference}
% Brief summary of Paper 1 findings. ~1 page.
% Three failures: observational data, dense prediction, VI collapse.
% GCV-ridge as the practical baseline this paper must beat.
\todo{Paper 1 summary}

\subsection{Instrumental Variables for Observational Causal Inference}
% Pearl's causal hierarchy. IV setup: E -> TF -> gene.
% Exclusion restriction. Instrument strength. 2SLS.
% Connection to condition-residualized representations.
% Honest caveats (Caveat 1, 2, 3 from brainstorm).
\todo{IV background}

\subsection{Perturb-seq Interventional Data}
% Replogle et al. 2022. K562 genome-wide CRISPRi.
% 9,866 knockdowns. 2.5M single cells. Pseudo-bulk construction.
% Delta y = y_k - y_0. Action labels a_k.
% Why this is P(Y|do(X)) and observational compendium is not.
\todo{Perturb-seq background}

\subsection{Sparsemax and Structured Sparsity}
% Sparsemax (Martins \& Astudillo 2016) vs softmax.
% Differentiable sparse probability simplex projection.
% Why sparsemax not top-K: differentiability, training stability.
% Convergence of nonzero support during training.
\todo{Sparsemax background}

% ============================================================
\section{Architecture}

\subsection{Design Principles}
% The five principles from the brainstorm document.
% Each principle is a direct response to a diagnosed failure from Paper 1.
% Table: failure -> principle -> architectural component.
\todo{Design principles}

\subsection{Stage 1 — Causal Context Encoder}
% Transformer encoder over D TF tokens.
% Token = TF identity embedding + expression value.
% Self-attention among TFs: learns regulatory crosstalk.
% Action conditioning: perturbed TF receives action embedding.
% Condition residualization: h_tf,d^resid = h_tf,d - Proj_condition(h_tf,d).
%   PCA of TF expression for observational IV version.
%   Explicit action label for interventional version.
% Output: H_tf in R^{D x d}.
% Why no EMA: direct supervised signal replaces self-distillation.
\todo{Stage 1 architecture}

\subsection{Stage 2 — Sparse Causal Predictor}
% For each target gene g:
%   Raw scores: e_{dg} = w_g^T tanh(W_d h_{tf,d}^resid)
%   Sparse attention: alpha_{dg} = sparsemax(e_{dg})
%   Context: context_g = sum_d alpha_{dg} h_{tf,d}
%   Prediction: delta_hat_g = context_g W_out + b_g   (expression space)
%
% Training objective:
%   L = sum_k sum_g (delta_hat_g^k - delta_y_g^k)^2  +  lambda * ||alpha||_0
%   MSE in expression space + sparsity penalty on attention support.
%
% Why expression space not latent: direct interpretability.
% Why sparsemax: differentiable sparse selection.
% Edge scores: s_{dg} = alpha_{dg} (primary) or |w_{dg}| * alpha_{dg} (refined).
\todo{Stage 2 architecture}

\subsection{Stage 3 — Bayesian Uncertainty Layer}
% NUTS on sparse support S_g = {d : alpha_{dg} > 0}.
% Per-gene regression: y_g = X_tf[:,S_g] beta_g + epsilon_g.
% Regularised horseshoe prior (Piironen \& Vehtari 2017).
%   Principled tau_0g per gene: (p_0 / (K_g - p_0)) * (sigma_g / sqrt(n)).
% Warm start: GCV-ridge point estimate from Paper 1.
%   GCV solution is near posterior mode -> short burn-in.
% NUTS settings: 4 chains, 1000 samples, target_accept=0.9.
%   Non-centred parameterisation for tau and lambda.
% Tractability argument: K_g * G vs D * G parameters.
%   K_g ~ 15, G ~ 4000: 60K vs 1.2M. NUTS feasible at 60K.
%
% Outputs per edge:
%   Posterior mean E[beta_{dg} | data]
%   95\% credible interval
%   Posterior inclusion P(|beta_{dg}| > threshold | data)
%   Variance Var(beta_{dg} | data) -- experimental design signal
\todo{Stage 3 architecture}

\subsection{Algorithm Summary}
% Algorithm box: full Causal B-JEPA training + inference pipeline.
\todo{Algorithm box}

% ============================================================
\section{Experimental Design}

\subsection{Two Training Regimes}
% Version A — Observational IV (DREAM5):
%   Training data: DREAM5 net3 E. coli expression compendium.
%   Causal identification: condition-residualized representations.
%   Instrument: experimental conditions as natural experiments.
%   Evaluation: DREAM5 gold standard AUPR + AUROC.
%
% Version B — Interventional (Replogle K562):
%   Training data: Replogle et al. 2022 K562_gwps CRISPRi.
%   Causal identification: true P(Y|do(X)).
%   Action conditioning: knockdown identity.
%   Evaluation: CausalBench + held-out knockdown prediction.
%
% Comparison: Version A vs Version B. Divergence = IV assumption violation.
\todo{Two versions}

\subsection{Baselines}
% GENIE3 (Paper 1 baseline).
% GCV-ridge (Paper 1 positive contribution -- the bar to beat).
% B-JEPA (Paper 1 primary method -- should be beaten substantially).
% SCENIC+ (state of the art for human GRN inference).
% scGPT regulatory head (if applicable).
\todo{Baselines}

\subsection{Evaluation Metrics}
% Primary (CausalBench):
%   Mean Wasserstein distance on held-out knockdowns.
%   False omission rate (FOR) on absent edges.
%
% Secondary (DREAM5):
%   AUPR. AUROC. Rank correlation with GCV-ridge.
%
% Novel (calibration):
%   95\% credible interval coverage on held-out Replogle knockdowns.
%   Target: empirical coverage = 0.95 if calibrated.
%
% Novel (experimental design):
%   Precision at top-k high-uncertainty moderate-confidence edges.
%   These are the novel interaction candidates for experimental validation.
\todo{Evaluation metrics}

% ============================================================
\section{Results}

\subsection{Version A: Observational IV on DREAM5}
% AUPR vs GCV-ridge, GENIE3, B-JEPA.
% Does condition residualization help vs no residualization?
% Sensitivity analysis: number of PCs removed.
% Does sparsemax attention find biologically plausible regulators?
% Figure: attention weight distribution. Top regulators per gene.
\todo{Version A results}

\subsection{Version B: Interventional on Replogle K562}
% CausalBench Wasserstein distance.
% Held-out knockdown prediction AUPR.
% Calibration coverage: does 95\% CI cover 95\% of true effects?
% Figure: calibration plot (empirical vs nominal coverage).
\todo{Version B results}

\subsection{Comparing the Two Versions}
% Where do A and B agree? These are causally robust edges.
% Where do they diverge? These are IV assumption violations.
% Biological interpretation of divergence points.
% Do divergence points correspond to known indirect regulatory effects?
\todo{Version comparison}

\subsection{Uncertainty-Guided Experimental Prioritization}
% Top-k edges by Var(beta_{dg}).
% High variance + moderate confidence = most informative to validate.
% Comparison against random selection and confidence-ranked selection.
% If wet lab validation available: precision on experimentally confirmed edges.
\todo{UQ results}

\subsection{Calibration Analysis}
% Figure: calibration coverage curve.
% 95\% CI coverage at different confidence thresholds.
% Comparison: NUTS vs VI vs bootstrap.
% Key claim: NUTS is calibrated, VI is not.
\todo{Calibration results}

% ============================================================
\section{Discussion}

\subsection{What Makes Causal B-JEPA Possible}
% Three enabling components that had to be resolved in Paper 1:
%   GCV-ridge warm start (from Paper 1 analytical horseshoe).
%   Failure analysis justifying sparse expression-space output.
%   IV framing justifying observational causal claims.
% Paper 1 is not just motivation -- it is the foundation.
\todo{What makes it possible}

\subsection{The Causal Identification Gap}
% Quantification of the A vs B difference.
% When is observational IV sufficient? When is interventional required?
% Practical guidance: if n/D > 3 and conditions are diverse, IV suffices.
% If edges involve indirect regulation or TF post-translational control, need Perturb-seq.
\todo{Causal identification gap}

\subsection{Limitations}
% Assumption 1: TF expression proxies TF regulatory activity (weakest for eukaryotes).
% Assumption 2: Exclusion restriction (approximate, adversarial with instrument strength).
% Assumption 3: Sparse regulatory structure (K_g ~ 15 -- may be wrong for master regulators).
% K562 -> E. coli transfer: not the claim. Version A is evaluated on DREAM5 directly.
\todo{Limitations}

\subsection{Connection to the Virtual Cell Program}
% CZI virtual cell initiative. Perturb-seq as training data.
% Causal B-JEPA as a module for regulatory structure discovery.
% Uncertainty quantification as the key gap in current virtual cell models.
% Future: multi-organism training, cell-type-specific GRN inference.
\todo{Virtual cell connection}

% ============================================================
\section{Conclusion}
\todo{Conclusion}

% ============================================================
\section*{References}
\bibliographystyle{abbrvnat}
\bibliography{references}
% Key citations to add:
% Paper 1 (this work -- cite arXiv preprint)
% Replogle et al. 2022 (Perturb-seq K562)
% Martins \& Astudillo 2016 (sparsemax)
% Piironen \& Vehtari 2017 (regularised horseshoe)
% Pearl 2009 (causality)
% Angrist \& Pischke 2009 (IV)
% CausalBench (Chevalley et al. 2023)
% SCENIC+ (Bravo Gonzalez-Blas et al. 2023)
% scGPT (Wang et al. 2024)
% LeCun 2022 (world models / JEPA)
% Huang 2026 (BiJEPA)

% ============================================================
\appendix

\section{Proof: Sparsemax Attention Converges to 2SLS Under Linear Conditions}
% Informal theorem: under linear encoder, exact condition residualization,
% and linear predictor, sparsemax attention weights converge to 2SLS coefficients.
% This connects the architecture to classical IV estimation theory.
\todo{2SLS convergence theorem}

\section{Replogle Dataset Details}
% K562 cell line. CRISPRi mechanism. Pseudo-bulk construction.
% Train/validation/test split. Which knockdowns are held out.
% Data preprocessing: normalisation, pseudo-bulk aggregation.
\todo{Replogle details}

\section{Hyperparameter Sensitivity}
% Sparsemax temperature. Number of PCA components removed.
% K_g prior specification. NUTS settings sensitivity.
\todo{Hyperparameter sensitivity}

\section{Full CausalBench Results}
% All metrics across all methods. Statistical significance.
\todo{Full CausalBench table}

\section{NUTS Diagnostics for Stage 3}
% R-hat, ESS, trace plots.
% Comparison: Stage 3 NUTS vs Paper 1 NUTS (warm start reduces burn-in by X\%).
\todo{NUTS diagnostics}

\end{document}
